\subsection{Perfect closures of global and henselian global fields}
\label{C:perfection-subsection}

The perfect-field results of
Theorems~\ref{C:perfect_Galois_indices} and
\ref{C:artinian_perfect_WG_reduction} retain genuine division-algebra
and hermitian arithmetic.  The following family supplies terminating
algorithms for those retained data.  Its coefficient fields are
usually neither global fields nor discretely valued fields.

\begin{theorem}[Proposed: effective arithmetic after perfection]
\label{C:perfection-field-arithmetic}
Let $K$ be either a supplied finite extension of an odd-characteristic
global function field, with the access of
Lemma~\ref{C:global-function-field-arithmetic}, or a supplied finite
extension of a henselian global field from
Theorem~\ref{C:henselian-global-field-arithmetic}.
Write $p=\operatorname{char}K$ and present
\[
 P=K^{\mathrm{perf}}=\bigcup_{n\ge0}K_n,
 \qquad K_n=K^{1/p^n},
 \qquad (a,n)\text{ names }a^{1/p^n}.                 \tag{C.3.PERF1}
\]
There are exact decidable finite names, arithmetic, and inverse
Frobenius on $P$, and ordinary terminating factorization algorithms
over every supplied finite field tower over $P$.  Validity of such
towers is decidable.

Over each such finite field, every finite central simple algebra has
computable index, an actual central division factor, matrix units,
and mutually inverse matrix-algebra maps.  Reduced norms and reduced
characteristic polynomials are computable.  Every finite-dimensional
algebra has a computable radical, a certified semisimple quotient with
actual division factors and centers, and actual lifted primitive
idempotents.  For each finite-dimensional division algebra with an
involution fixing $P$, nonsingular hermitian isometry of either sign
is decidable and every positive answer constructs an actual isometry.
The least-Gram and least-pair arithmetic of
Theorem~\ref{C:henselian-hermitian-isometry} is consequently ordinary
and terminating over these division algebras.
\end{theorem}
\begin{proof}
\emph{Names and finite stages.}
The map $K\longrightarrow K_n$, $a\longmapsto a^{1/p^n}$, is a field
isomorphism.  At any fixed stage all the algorithms already proved
for $K$ and its finite extensions therefore apply by relabeling
coefficients.  If $n\le m$, the transition sends the name $a$ in
$K_n$ to the name $a^{p^{m-n}}$ in $K_m$.  To compare two names or
perform an operation, move them to a common stage and use exact
arithmetic in $K$.  Inverse Frobenius sends $(a,n)$ to $(a,n+1)$.
These are finite operations and give an enumeration of $P$ with
decidable equality and valid names.

We describe a finite-stage property that will be used repeatedly.
Let $L/P$ be an already valid finite field tower.  Since $P$ is
perfect, its nontrivial defining irreducible polynomials are separable.
Put all their coefficients in one sufficiently large $K_n$ and use
the same tower equations there, obtaining $L_n/K_n$.  At each tower
step the polynomial remains irreducible over the smaller field:
a factorization there would still be a factorization over the
corresponding $P$-field.  Its nonzero derivative is unchanged.
Thus $L_n/K_n$ is a valid finite separable tower, with
\[
 L=L_nP,\qquad L_m=L_nK_m\quad(m\ge n),\qquad
 [L_m:K_m]=[L_n:K_n].                                  \tag{C.3.PERF2}
\]
The degree equality follows from linear disjointness of a finite
separable extension and a purely inseparable extension.  Hence these
are actual fields and their tower bases are compatible.  The union
of the $L_m$ is $L$, and $L/L_n$ is purely inseparable.  A finite
list of elements of $L$ is put in a common $L_m$ simply by inspecting
its coefficient names in that tower basis.  Every finite $L_n$-field
inside this union has $p$-power degree.

\emph{Factorization without a circular tower test.}
Assume first that the field $L$ is already valid; this includes the
base case $L=P$.  For a polynomial $f\in L[X]$, handle zero and
constant inputs separately and divide a nonconstant input by its
leading coefficient.  Put its coefficients and the tower in a common
stage $L_n$.  The previously proved finite-extension arithmetic over
$K_n\simeq K$ factors $f$ there with its multiplicities.

For a monic irreducible factor $g\in L_n[X]$, write
$g(X)=h(X^{p^e})$ with $e\ge0$ maximal.  The polynomial $h$ is
separable and irreducible over $L_n$.  A separable irreducible
polynomial remains irreducible over any purely inseparable extension,
so $h$ is irreducible over $L$.
The field $L$, being a finite extension of a perfect field, is perfect.
Its inverse Frobenius is computable: enumerate its exact elements
until $y^p=x$; existence and uniqueness prove termination.
Apply this operation $e$ times to every coefficient of $h$ and call
the resulting monic polynomial $\widetilde h$.  Then
\[
                         g=\widetilde h^{\,p^e}.         \tag{C.3.PERF3}
\]
The polynomial $\widetilde h$ is irreducible over $L$: applying the
Frobenius automorphism to a hypothetical proper factorization would
factor $h$.  This gives all irreducible factors, their actual
coefficients, and their multiplicities.  Equal factors are collected
by exact equality, and the original leading coefficient is restored.
Degree-one inputs are immediate.

To validate a proposed tower, begin with the established field $P$,
factor its first defining polynomial, and admit the new field only
if that polynomial is irreducible.  Then repeat over the field just
validated.  Its perfectness is a mathematical consequence of finite
extension, and its arithmetic is its supplied finite vector-space
table.  Thus the inverse-Frobenius search above is used only over
an already valid perfect field.  The finite-stage construction
(C.3.PERF2) is likewise used only after validity has been established.
This proves the factorization and tower assertions inductively.

\emph{A terminating test for the division index.}
Let $A$ be a central simple $L$-algebra of degree $N$.
Put all its structure constants in a common $L_n$ and obtain $A_n$.
Associativity and the unit identities descend by exact equality.
The invertibility of
$A\otimes_L A^{\mathrm{op}}\longrightarrow\operatorname{End}_L(A)$
is a nonzero finite determinant; the same determinant is nonzero
over $L_n$.  Thus $A_n$ is central simple.  For $m\ge n$ put
$A_m=A_n\otimes_{L_n}L_m$ and compute $d_m=\operatorname{ind}(A_m)$
using the established global or henselian-global central-simple
algorithms over $L_m/K_m$.

We prove that the following stopping rule is valid and eventually
occurs:
\[
              \text{stop at the first }m\ge n
                         \text{ with }p\nmid d_m.       \tag{C.3.PERF4}
\]
For any finite extension $E/F$ and degree-$N$ central simple algebra,
write $d_F,d_E$ for its two indices.  Base extension of a division
factor shows $d_E\mid d_F$.  A simple module over the extended
algebra has $E$-dimension $Nd_E$; viewed over $F$, it is a module
for the original simple algebra and has dimension $[E:F]Nd_E$.
Every such module has $F$-dimension divisible by $Nd_F$.
Consequently
\[
                   d_E\mid d_F\mid [E:F]d_E.            \tag{C.3.PERF5}
\]
In particular a finite purely inseparable extension changes an
index only by a power of $p$.
The same assertion about the prime-to-$p$ part holds for arbitrary
purely inseparable extensions: complete matrix units displaying the
new division factor use only finitely many coefficients and therefore
descend to a finite purely inseparable intermediate field.  The
descended corner is already division, since nonzero zero divisors
would remain zero divisors after extension.  Apply (C.3.PERF5) there.

The field $L$ is perfect.  Its separable closure is then perfect,
so $x\mapsto x^p$ is an automorphism of its multiplicative Galois
module.  Via the cohomological Brauer identification this implies
that multiplication by $p$ on $\operatorname{Br}(L)$ is injective;
there is no nonzero $p$-primary Brauer class.  The period and index
have the same prime divisors.  These precise facts follow from
\cite[Theorem 4.4.7 and Proposition 4.5.13, pp.99,104--105]
{C-GilleSzamuelyPerfection2006}; the displayed Frobenius argument
also proves the needed perfect-field consequence directly.
It follows that $p\nmid d_\infty$, where
$d_\infty=\operatorname{ind}(A)$.

Choose complete matrix units over $L$ for
$A\simeq M_{N/d_\infty}(D_\infty)$.  They and their finite
coordinate inverses lie in some $A_m$.  Their corner there has
degree $d_\infty$ and is division by the preceding reflection
argument.  Hence $d_m=d_\infty$, so (C.3.PERF4) terminates.
Conversely, whenever $p\nmid d_m$, invariance of the prime-to-$p$
part forces $d_m=d_\infty$.  The actual division factor and matrix
units computed at this stopping stage therefore remain a division
factor and matrix units after extension to $L$.
The same argument proves the useful formula
$d_\infty=d_n/p^{v_p(d_n)}$, but the procedure still finds an
actual finite stage and actual matrices, rather than merely an index.

Finite separable splitting-field certificates, now using the proved
tower factorization, compute the reduced norm and characteristic
polynomial exactly as in
Theorem~\ref{C:global-function-central-simple-algorithms}.
Every certificate is a finite multiplication-table isomorphism and
its inverse.  The determinant in a splitting matrix algebra descends
to the center; equality and its center coordinates are exact.

\emph{Finite algebras and actual radicals.}
Return the zero algebra, its zero radical and its empty factor list
separately.
For a finite-dimensional $L$-algebra $E$, enumerate finite center
field towers $T_j/L$, central simple $T_j$-algebras $C_j$ certified
by the invertible multiplication map, and unital multiplicative
surjections $E\to\prod_j C_j$.  Require that the kernel have zero
$(\dim_LE+1)$st power.  All tests are finite exact linear algebra
and the now established field-tower tests.  A successful kernel is
nilpotent and the quotient is semisimple, so the kernel is exactly
the Jacobson radical.  Artin--Wedderburn gives a successful finite
certificate; no semisimple subalgebra of $E$ is presumed.
Apply (C.3.PERF4) to each $C_j$ to obtain its actual division
factor and matrix units.  Primitive idempotents lift through the
nilpotent radical, and exact positive enumeration, or the earlier
nilpotent lifting formula, computes those lifts.  This is the full
certificate procedure of
Theorem~\ref{C:global-function-finite-algebra-algorithms}, with
every arithmetic gate now supplied for $L$.

\emph{Hermitian isometry reflects to one finite stage.}
Let $D$ be a supplied finite-dimensional division algebra over $P$,
with involution $\sigma$ fixing $P$, and let $h,g$ be nonsingular
$\varepsilon$-hermitian forms.  Its center is a finite extension
of $P$.  Put a $P$-basis multiplication table for $D$, the table
of $\sigma$, and all Gram entries in one $K_n$; call the descended
algebra with involution $(D_n,\sigma_n)$ and the forms $h_n,g_n$.
The map $D_n\to D$ is injective and $D_n\otimes_{K_n}P=D$.
A nonzero element of $D_n$ with singular left-multiplication
matrix would still have singular left multiplication over $P$,
contradicting division.  Thus $D_n$ is already a division algebra.
Likewise the forms are nonsingular at this stage.  Its center and
the involution-fixed subfield of that center are computed by linear
kernels; they are finite fields over $K_n$ and receive actual tower
presentations by successive basis adjunction.

The already proved global or henselian-global hermitian algorithms
therefore decide $h_n\simeq g_n$.  If they are isometric, the
computed isometry extends to $P$.  Conversely an isometry over $P$
has finitely many scalar coefficients and hence is defined over a
finite purely inseparable extension $M/K_n$ inside $P$.  Its degree
is a power of the odd prime $p$.  Bayer--Fluckiger--Lenstra's
odd-degree descent theorem gives
\[
 h_n\otimes M\simeq g_n\otimes M
                 \quad\Longrightarrow\quad h_n\simeq g_n.
                                                               \tag{C.3.PERF6}
\]
Their Proposition 1.2 and Corollary 1.4 apply to every finite
odd-degree extension, including purely inseparable ones, and to
either hermitian sign over any finite-dimensional division algebra
with base-linear involution
\cite[pp.362--363]{C-BayerLenstraOddDegree1990}.
The proof uses a nonzero linear transfer, not the field trace, so
no separability has been inserted.  Formula (C.3.PERF6) proves that
this single finite-stage decision is an exact decision over $P$,
with an actual positive isometry.  The same argument applies to a
division algebra over any finite $P$-field whose involution fixes
$P$, by viewing its table over $P$ as above.

Finally, all Gram conditions, ranks, signs and finite codes are
decidable.  Testing the finitely many preceding valid codes gives
least Gram representatives.  Hermitian cancellation over a division
algebra turns equality of pairs into isometry of their cross sums,
so the least-pair groups and their operations are obtained exactly
as in Theorem~\ref{C:henselian-hermitian-isometry}.  Unequal ranks
and common rank zero are handled before the positive-rank test.
This proves every asserted algorithm.
\end{proof}

\begin{theorem}[Proposed: complete split groups over perfect-closure coefficients]
\label{C:perfection-coefficient-WG}
Let $P$ be as in Theorem~\ref{C:perfection-field-arithmetic}, and
let $k/P$ be a supplied finite field tower.  For every $m\ge0$,
complete ordinary terminating classification and split-$WG$
operation algorithms for rational $O_{m,R}$-representations hold
for the following coefficient rings:
\begin{enumerate}
\item every finite-dimensional commutative $k$-algebra given by
its multiplication table;
\item $R=B^h$, where $B$ is a countable commutative local ring
with exact decidable finite names, arithmetic, residue map and
an ordinary algorithm giving finite generators of every finite
matrix kernel, and the residue field is a supplied finite
extension of $k$.
\end{enumerate}
The output consists of a decidable literal index set
$\mathscr I_R\subseteq\mathbf N$ of actual projective
indecomposables, their dual indices, signs and reference forms,
actual residual division algebras $D_i$ and adjoints $\sigma_i$,
and the additive coordinates
\[
 WG_{\rm sp}(BO_{m,R})\simeq
 \bigoplus_{i=i^*}\mathcal G(D_i,\sigma_i,\varepsilon_i)
 \ \oplus\
 \bigoplus_{\{i,i^*\},\,i\ne i^*}\mathbf Z.             \tag{C.3.PERF7}
\]
Here $\mathcal G$ is the literal least-Gram-pair group, now with
the decisions of Theorem~\ref{C:perfection-field-arithmetic}.
Actual equivariant isometry, equality of finite virtual inputs,
and every specified product or exterior coefficient are decidable
and computable.  Every finite coordinate vector has actual positive
and negative representatives.
\end{theorem}
\begin{proof}
We verify the inputs to the full construction of
Theorem~\ref{C:henselian-global-coefficient-WG}; its earlier
discrete-valuation hypothesis is replaced here by the preceding
arithmetic theorem, not asserted for $P$.

For a finite coefficient algebra, present modules by finite
$k$-spaces with the coefficient-algebra action and finite-degree
rational coaction matrices.  All identities are finite.  Finite
projectivity is decided by the affine linear equations for a
section of the canonical surjection $R^s\to M$ obtained from a
$k$-basis.  Compute actual Hom spaces, forms, and the $R$-dual by
finite linear kernels.  On a matrix space of size $s$, a grid of
$s+1$ distinct elements of the infinite field $k$ decides whether
its determinant is identically zero, and constructs a unit whenever
one exists.  Thus module isomorphism and both nonsingular form
signs are decidable by actual intertwiners.

For $B^h$, the pointed standard-etale exact-access proof of
Theorem~\ref{C:effective-coherent-henselization-WG} uses the given
finite matrix kernels, finite free monic algebras, residue arithmetic,
and faithful flatness, without a $C_1$ assumption.  It therefore
gives exact names and every full finite matrix kernel over $R$ here.
Reduce the computed actual Hom generators to the residue field
$\kappa$, which is finite over $P$.  The determinant grid in their
actual image decides units and sign forms, and coefficient lifting
gives actual matrices.  One does not substitute all intertwiners
of the reduced representations for this actual image.

The finite-dimensional endomorphism algebra in the first case has
its complete radical and actual semisimple factors by
Theorem~\ref{C:perfection-field-arithmetic}.  In the local case, if
$E_0$ is the actual residue image of the finite endomorphism ring
$E$, its kernel is radical: $1-au$ has unit determinant for every
kernel element $u$ and every $a\in E$, and its inverse remains
equivariant.  Thus $E/\operatorname{rad}E=E_0/\operatorname{rad}E_0$.
The same theorem computes this quotient over $\kappa$.
Lemma~\ref{C:nonnoetherian-henselian-lifting} supplies actual
idempotent lifts and semiperfectness, without a Noetherian condition.
Enumerate actual endomorphisms with the specified residue until an
idempotent lift is found.  The proved existence makes this a
terminating search.  Lifted idempotents split actual projective
submodules; positive dimension or local rank decreases at each split.
This computes the full Krull--Schmidt decomposition.

Finite predecessor isomorphism tests give the literal indecomposable
codes and their duals.  For a self-dual indecomposable, the two
forms $g+g^{\mathsf t}$ and $g-g^{\mathsf t}$ associated to an
isomorphism with its dual, after composition with $g^{-1}$, cannot
both be nonunits in its local endomorphism ring, because their sum
is $2$.  The determinant tests select a sign and an actual reference
form; finite predecessor tests bounded by this witness select $q_i$.
Its adjoint gives the explicit involution on the actual division
quotient $D_i$; this involution fixes $k$, or $\kappa$ in the local
case, and in particular fixes $P$.
The last part of Theorem~\ref{C:perfection-field-arithmetic}
therefore decides every residual multiplicity-form isometry and
constructs its actual matrix.

The checked hermitian transfer and radical-lifting proof now gives
exactly (C.3.PERF7).  For clarity, lifting a residual signed Gram
matrix $H$ means lifting its entries to the actual endomorphism
ring and replacing a lift $H_0$ by
$(H_0+\varepsilon_i H_0^\dagger)/2$.  Its residue is the prescribed
invertible $H$, so the lift is invertible.  Multiplying by the
reference form yields the actual symmetric representation.
The radical isometry-lifting theorem identifies precisely the
same residual forms.  For a positive isometry decision one may
also enumerate all actual equivariant matrices until the exact
Gram equation is met; the lifting theorem proves termination.
Non-self-dual pairs contribute the actual hyperbolic object.
Cancellation in each residual division block and uniqueness by
dual orbit give the stated least-pair equality decision and actual
positive and negative representatives.

All operations are taken on these actual representations first.
Balanced tensor products and exterior powers over $R$ have finite
presentations and actual coactions.  The invariant forms are the
tensor Gram matrices and the determinant Gram matrices on exterior
powers; their nonsingularity follows locally for finite projectives.
No division by an exterior factorial is used.  Decompose the result
again, reduce the actual multiplicity matrices in the computed
division quotients, and canonicalize by the proved isometry tests.
For virtual exterior operations, the coefficient of a specified
degree in $\lambda_t(x)\lambda_t(y)^{-1}$ uses only finitely many
coefficients and the finite inverse-series recurrence.
Zero rings and zero modules are returned separately as in the
earlier construction.  Every step has now been supplied by an
ordinary terminating algorithm in the present coefficient family.
\end{proof}

\begin{corollary}[Proposed: perfect non-$C_1$ fields of nondiscrete value group]
\label{C:perfect-global-nonC1-family}
For each odd prime power $q$, both
\[
 P_q=\mathbf F_q(t)^{\mathrm{perf}},\qquad
 Q_q=\left(\operatorname{Frac}
                  (\mathbf F_q[t]_{(t)}^h)\right)^{\mathrm{perf}}
                                                               \tag{C.3.PERF8}
\]
are concrete new perfect non-$C_1$ coefficient fields for the
complete algorithms above, together with all their supplied finite
field extensions, finite commutative coefficient algebras, and the
effective coherent-local henselizations specified in the theorem.
The natural $t$-valuations have residue field $\mathbf F_q$ and
value group $\mathbf Z[1/p]$.  The latter field is henselian for
this valuation; the former is not.  Neither is a global function
field or a finite extension of one of the previously used henselian
global fields.
\end{corollary}
\begin{proof}
The arithmetic assertions follow directly from the two choices of
$K$ in Theorem~\ref{C:perfection-field-arithmetic}.
On a named element $a^{1/p^n}$ the extended valuation is
$v(a)/p^n$.  Each field contains all $t^{1/p^n}$, so its value group
is exactly $\mathbf Z[1/p]$.  Its residue is still $\mathbf F_q$:
a purely inseparable residue extension of the perfect field
$\mathbf F_q$ is trivial.  Algebraic extensions of a henselian
valued field are henselian, proving this assertion for $Q_q$.

Choose a nonsquare $a\in\mathbf F_q^\times$.  Over either field,
for any $(x,y)\ne(0,0)$ the binary norm has valuation
\[
                       v(x^2-ay^2)=2\min(v(x),v(y)).    \tag{C.3.PERF9}
\]
Indeed after division by the square of a coordinate of least
valuation, the two residues are not both zero, and the residual
binary form is anisotropic over $\mathbf F_q$.
Consequently
\[
                       X_1^2-aX_2^2-tX_3^2+atX_4^2    \tag{C.3.PERF10}
\]
has no nonzero zero: when both binary pairs are nonzero, equality
would place the same valuation in both
$2\mathbf Z[1/p]$ and $1+2\mathbf Z[1/p]$.
These are disjoint since $p$ is odd.  If one binary pair is zero,
(C.3.PERF9) forces the other to be zero as well.
This is an anisotropic quadratic form in four variables, proving
that both fields are not $C_1$.  In particular no disappearing
quasilinear $p$-form is used as a witness after perfection.

The polynomial $X^2-(1+t)$ has the simple root $1$ modulo the
maximal ideal of the $t$-valuation on $P_q$, but no root in $P_q$.
For if $b^2=1+t$ with $b^{p^n}\in\mathbf F_q(t)$, then
$(1+t)^{p^n}$ would be a square in $\mathbf F_q(t)$, whereas its
valuation at the prime $1+t$ is the odd integer $p^n$.
This proves the asserted nonhenselian property of $P_q$.

Finally a perfect positive-characteristic field admits no nontrivial
discrete valuation: its value group is $p$-divisible because every
element has a $p$th root.  Global function fields and the earlier
finite henselian-global fields admit nontrivial discrete valuations,
so neither field in (C.3.PERF8) is such a field.  Equivalently,
the earlier fields have positive imperfection degree.
Characteristic excludes the rational-algebra family, and
(C.3.PERF10) excludes every $C_1$ family.  The earlier general
perfect-field Galois and scalar presentations retained the arithmetic
gates now resolved by Theorem~\ref{C:perfection-field-arithmetic}.

In particular the nonfields $P_q[\epsilon]/(\epsilon^s)$ and
$Q_q[\epsilon]/(\epsilon^s)$ for $s\ge2$, as well as local
henselizations such as $P_q[z]_{(z)}^h$ and $Q_q[z]_{(z)}^h$,
have the full output.  More generally, the finite kernel algorithms
for finitely presented affine algebras over these effective fields
and pointed localization give the stated affine local examples,
including singular and nonreduced ones whenever the supplied
residue field is finite over the indicated coefficient field.
All input records and each requested operation are finite; no
algorithm for infinitely supplied products is asserted.
\end{proof}
